\documentclass[11pt]{article}

\usepackage[margin=0.9in]{geometry}
\usepackage{amsmath}
\usepackage{booktabs}
\usepackage{graphicx}
\usepackage{float}
\usepackage[hidelinks]{hyperref}
\usepackage{array}
\usepackage[table]{xcolor}
\usepackage{parskip}
\usepackage{titlesec}

\graphicspath{{../../figures/}}
\sloppy
\setcounter{tocdepth}{2}

\newcommand{\mainfigure}[3]{%
\begin{figure}[H]
    \centering
    \includegraphics[width=\textwidth]{#1}
    \caption{#2}
    \label{#3}
\end{figure}
}


\title{FCFS Appointment Simulation:\\
Metric-Focused Sensitivity Analysis\\[0.4em]
\large Supervisor Review Version}
\author{CUIMC Appointment Simulation}
\date{\today}

\begin{document}
\maketitle

\begin{abstract}
\noindent
A first-come, first-served (FCFS) outpatient appointment simulation with two
symmetric patient classes was run under a wide range of behavioral assumptions.
The central finding is that utilization and access are distinct outcomes: at
baseline, the system uses 83.9\% of available slots but serves only 26.9\% of
arriving patients. Total demand pressure is the strongest driver of overall
served rate, while class-specific behavioral differences---primarily
cancellation probability, balking threshold, and no-show step---are the
dominant drivers of class disparity. This memo explains the metric hierarchy,
reports the main sensitivity findings, and identifies the model's key
limitations.
\end{abstract}

\tableofcontents
\newpage

%% ============================================================
\section{Research Question and Interpretation Framework}
%% ============================================================

\subsection*{Research Question}

How does FCFS appointment scheduling perform across a range of demand and
behavioral assumptions, and which patient-level behavioral parameters most
strongly affect access and equity?

\subsection*{Why Multiple Metrics Are Needed}

A single metric cannot characterize appointment system performance because
utilization, patient access, patient-facing wait, and class equity are
logically independent outcomes. A clinic can be fully booked (high
utilization), yet serve a small fraction of requesting patients (low access),
if many patients balk at long offered waits or cancel after booking. Similarly,
average wait times can appear short if patients with long offered delays are
systematically excluded from the accepted-booking pool.

The simulation therefore tracks five conceptually distinct outcomes:
\begin{enumerate}
    \item \textbf{Capacity use}: how busy is the clinic?
    \item \textbf{Patient access}: what share of patients who want service
    actually receive it?
    \item \textbf{Patient-facing wait}: how long must an average arriving
    patient wait?
    \item \textbf{Rejection behavior}: how often do patients reject offered
    appointments, and why?
    \item \textbf{Class equity}: do different patient groups experience
    systematically different access?
\end{enumerate}

Each metric is necessary to interpret the others, and none is sufficient
on its own.

\subsection*{Metric Hierarchy}

\begin{table}[H]
\centering
\renewcommand{\arraystretch}{1.2}
\begin{tabular}{p{0.22\textwidth}p{0.22\textwidth}p{0.25\textwidth}p{0.22\textwidth}}
\toprule
Metric & What it measures & How to interpret it & Main caveat \\
\midrule
Overall served rate & Share of arriving patients who complete service & Primary access metric; lower is worse & Denominator is all arrivals, not just those who book \\
Average utilization & Share of clinic slots filled with completed visits & Capacity-use metric; high does not imply good access & No-shows and balking inflate utilization relative to access \\
Mean offered wait & Average delay offered to all patients who receive a slot offer & Patient-facing wait, inclusive of those who then balk & Harder to interpret when demand mix changes \\
Mean accepted wait & Average delay among patients who accepted a booking & Conditional delay metric & Selection bias: excludes patients who reject long offers \\
Balking rate & Share of offered patients who reject the offered slot & Congestion and rejection diagnostic & Not a success metric; lower is not automatically better \\
Class served-rate gap & Class 1 served rate minus Class 2 served rate & Fairness / class disparity metric & Meaningful only when class behavioral parameters differ \\
\bottomrule
\end{tabular}
\caption{Metric hierarchy. Primary access metric is overall served rate; all
others contextualize it.}
\label{tab:metric-hierarchy}
\end{table}

%% ============================================================
\section{Baseline Diagnosis}
%% ============================================================

\subsection{Simulation Setup}

The baseline uses \texttt{configs/baseline.yaml}: 32 slots per day, a 14-day
booking horizon, two symmetric classes each with 50 arrivals per day,
cancellation probability 0.10, balking threshold 9 days with high-delay
balking probability 0.50, and no-show threshold 6 days with high-delay
no-show probability 0.30. Low-delay balking and no-show probabilities are
both 0.00. The measurement window is 365 days following a 30-day burn-in.

\subsection{Baseline Metric Summary}

\begin{table}[H]
\centering
\renewcommand{\arraystretch}{1.15}
\begin{tabular}{lrrrrrr}
\toprule
Avg.\ utilization & Overall served & Mean accepted wait & Mean offered wait & Class gap & Delay gap \\
\midrule
0.839 & 0.269 & 8.35 & 9.30 & 0.001 & 0.003 \\
\bottomrule
\end{tabular}
\caption{Baseline metric outcomes. Class gap is Class~1 served rate
minus Class~2 served rate. Delay gap is Class~2 mean offered wait
minus Class~1 mean offered wait.}
\label{tab:baseline-outcomes}
\end{table}

\textbf{The key diagnostic is the utilization--access gap.} At baseline, the
system uses 83.9\% of available slots but serves only 26.9\% of arriving
patients. This gap indicates that the system is effectively capacity-constrained:
the booking horizon fills quickly, large numbers of patients receive no slot
offer or balk at long offered delays, and after-booking losses (cancellations
and no-shows) further reduce completed visits. A facility-manager who only
monitors utilization would conclude the clinic is functioning well; the
patient-side metrics tell a different story.


%% ============================================================
\section{Sensitivity Analysis}
\label{sec:sensitivity}
%% ============================================================

\subsection{Overview}

The sensitivity analysis varies one group of assumptions at a time around the
baseline to isolate how each driver moves each metric. All runs use the same
FCFS engine; only the input parameters and random seeds change. Each result
is averaged across multiple seeds to reduce sampling noise. The analysis is a
model-based stress test, not an empirical causal estimate.

\subsection{Average Utilization}

\subsubsection*{Definition}

\[
\texttt{average\_utilization} =
\frac{1}{D}\sum_d \frac{\text{served slots on day } d}{\text{slots per day}}.
\]

No-shows do not count as utilized slots because the visit was not completed.

\subsubsection*{Main Drivers and Interpretation}

Figure~\ref{fig:metric-util-drivers} shows that no-show behavior is the
clearest direct driver of utilization: booked slots that become no-shows
never convert to completed visits. Cancellation also matters because it
changes the appointment pool before the service day. Demand pressure, by
contrast, has a weaker effect on utilization once the system is already
operating near capacity---additional arrivals are simply turned away or balk,
leaving the filled slots largely unchanged.

The operational implication is that a provider concerned primarily about
utilization should focus on no-show reduction and cancellation management
rather than on reducing demand. However, high utilization achieved by filling
slots with patients who would have been lost to no-shows or cancellations does
not directly translate into better patient access.

\mainfigure{metric_utilization_drivers.png}{%
Utilization remains near 0.839 across a wide demand range, while
overall served rate falls sharply---confirming that utilization is a
poor proxy for patient access. Class 1 varies on the x-axis; Class 2
stays at baseline.}{fig:metric-util-drivers}

\subsection{Overall Served Rate}

\subsubsection*{Definition}

\[
\texttt{overall\_percent\_serviced} =
\frac{\sum_i V_i}{\sum_i A_i},
\]

where $A_i$ is arrivals and $V_i$ is completed visits for class $i$. This
is the primary access metric.

\subsubsection*{Main Drivers and Interpretation}

Figure~\ref{fig:metric-access-drivers} shows that total arrival pressure is
the dominant aggregate driver. As demand rises against fixed 32-slot daily
capacity and a 14-day booking horizon, an increasing share of arrivals either
receives no slot offer (the horizon is full) or balks at long offered delays.
No-shows and cancellations reduce completed visits after booking.

The regression screen (Section~\ref{sec:regression}) confirms this ordering:
total arrival rate has the largest standardized coefficient for overall served
rate ($\beta = -0.774$), followed by average no-show threshold ($\beta =
0.251$) and no-show step ($\beta = -0.175$).

Near the baseline, the system is already capacity-constrained: virtually all
slots are occupied yet fewer than 30\% of patients are served. This means
that incremental increases in demand primarily worsen unmet access rather
than further increasing utilization. Conversely, modest reductions in
arrival load---or increases in daily capacity---would have an outsized effect
on served rate in this regime.

\mainfigure{metric_access_drivers.png}{%
Overall served rate falls as total demand or class-specific arrival
rates rise, even when utilization stays high. At baseline, 50
arrivals per class per day produces a served rate near 0.269, well
below the utilization of 0.839.}{fig:metric-access-drivers}

\subsection{Mean Offered Wait}

\subsubsection*{Definition}

\[
\texttt{mean\_offered\_booking\_delay} =
\frac{\sum_i \text{total offered delay}_i}{\sum_i O_i},
\]

where $O_i = B_i + K_i$ is the number of offered patients: those who
accepted plus those who balked. Patients who received no slot offer are
excluded because no delay was presented to them.

\subsubsection*{Why Offered Wait Is Preferred Over Accepted Wait}

Mean accepted wait is computed only over patients who accepted a booking.
When balking rates are high, the accepted sample is systematically selected
against patients who received---and rejected---long-delay offers. Accepted
wait therefore understates the delay that arriving patients face.

Mean offered wait corrects for this selection by including balked patients.
It answers the question: among patients who received any delay offer, what
was the average delay presented? Figure~\ref{fig:metric-wait-drivers} shows
that offered wait is more sensitive to demand pressure than accepted wait,
because rising demand pushes more offers to the far end of the booking horizon
before patients reject them.

\mainfigure{metric_wait_drivers.png}{%
Mean offered wait rises with demand pressure and falls when balking
selects out long-delay patients. Mean accepted wait can fall even as
access worsens because patients who reject long-delay offers are
removed from the accepted sample.}{fig:metric-wait-drivers}

The regression screen finds total arrival rate ($\beta = 0.576$) as the
dominant driver of offered wait, followed by average cancellation probability
($\beta = -0.441$, note: cancellations free future capacity and reduce
congestion), average balking threshold ($\beta = 0.278$), and average balking
step ($\beta = -0.206$). The negative sign on balking step is a selection
effect: higher step means more long-delay patients reject offers, which
removes them from the offered-wait average.

\subsection{Balking Rate}

\subsubsection*{Definition}

\[
\texttt{balking\_rate} = \frac{\sum_i K_i}{\sum_i O_i}.
\]

Measured among patients who received an offer. Patients who received no offer
are excluded.

\subsubsection*{Interpretation as a Diagnostic}

Balking rate measures how often patients reject offered appointments. It is
a useful congestion diagnostic---rising balking rate indicates that offered
delays are routinely exceeding patient tolerance---but it is not a success
metric. A lower balking rate does not imply better access if it results from
fewer patients receiving offers, reduced arrival rates, or a selection shift
in who presents.

Figure~\ref{fig:metric-balk-rate-drivers} shows that balking rate is
controlled mainly by the balking step (probability increase above threshold)
and the threshold itself. The key interpretive warning is that lower
accepted wait often accompanies higher balking: when more long-delay patients
leave the system, the remaining accepted bookings are shorter on average.
This is a selection effect, not a service improvement.

\mainfigure{metric_balking_rate_drivers.png}{%
Balking rate rises with the high-balking probability and falls with
the threshold. Higher balking can coexist with lower accepted wait,
because rejecting long-delay patients removes them from the accepted
sample without reducing offered delays for remaining patients.
}{fig:metric-balk-rate-drivers}

\subsection{Class Served-Rate Gap}

\subsubsection*{Definition}

\[
\Delta_{\text{access}} =
\texttt{percent\_serviced}_1 - \texttt{percent\_serviced}_2.
\]

Positive values indicate Class 1 is served more often than Class 2.

\subsubsection*{When Class Labels Matter}

In the FCFS baseline, both classes have identical behavioral parameters.
Class labels are therefore exchangeable: the expected class gap is zero, and
observed deviations are noise. Class disparity becomes meaningful only when
the classes receive different behavioral assumptions.

Figure~\ref{fig:metric-class-gap-drivers} illustrates this directly: holding
Class 2 fixed at baseline while varying Class 1 assumptions, the class gap
widens as Class 1 behavioral parameters diverge from Class 2. Higher Class 1
cancellation probability, higher balking step, or higher no-show step all move
the gap against Class 1.

The regression screen (Section~\ref{sec:regression}) ranks the gap drivers:
cancellation probability gap ($\beta = -0.452$), balking threshold gap
($\beta = 0.429$), balking step gap ($\beta = -0.349$), no-show threshold
gap ($\beta = 0.268$), and no-show step gap ($\beta = -0.211$) are all
significant. Class~1 arrival share ($\beta = -0.120$) also enters
significantly, reflecting that FCFS is sensitive to the temporal ordering
of arrivals as well as their count.

\mainfigure{metric_class_gap_drivers.png}{%
The class served-rate gap is driven by differences in behavioral
parameters. With symmetric assumptions, both class lines track
together. As Class 1 behavioral parameters diverge from Class 2
(e.g., higher cancellation or balking), Class 1 access falls below
Class 2.}{fig:metric-class-gap-drivers}

%% ============================================================
\section{Balking Deep Dive}
%% ============================================================

This section isolates Class~1 balking behavior while Class~2 stays at
baseline. The purpose is to demonstrate the selection mechanism that makes
lower accepted wait a potentially misleading signal.

\subsection{Class 1 High Balking Probability}

\mainfigure{../../../../../outputs/class1_balking/figures/percent_serviced_by_class.png}{%
As Class 1 high-balking probability increases, Class 1 access falls
while Class 2 access rises. The redistribution shows that higher
balking worsens access for the more-balking class and partially
relieves congestion for the other.}{fig:balk-deep-step-served}

\mainfigure{../../../../../outputs/class1_balking/figures/mean_accepted_delay_by_class.png}{%
Class 1 accepted delay falls as high-balking probability increases,
but this is a selection effect: long-delay Class 1 patients leave the
accepted sample, not because slots become available sooner.
}{fig:balk-deep-step-accepted}

\mainfigure{../../../../../outputs/class1_balking/figures/mean_offered_delay_by_class.png}{%
Mean offered delay captures the broader congestion effect.
Once enough long-delay Class 1 patients reject offers, congestion
falls and offered delays shorten for both classes.
}{fig:balk-deep-step-offered}

The three figures together demonstrate that accepted wait and served rate
can move in opposite directions when balking changes. A policy intervention
that increases balking step would appear to improve Class~1 wait times while
actually worsening Class~1 access.

\subsection{Class 1 Balking Threshold}

\mainfigure{../../../../../outputs/class1_balking_threshold/figures/percent_serviced_by_class.png}{%
Raising Class 1's balking threshold increases Class 1 access by
tolerating longer offered waits. The benefit comes at the cost of
slightly longer congestion for Class 2 when the threshold is very
high.}{fig:balk-deep-threshold-served}

\mainfigure{../../../../../outputs/class1_balking_threshold/figures/service_gap.png}{%
The Class 1 minus Class 2 service gap is positive when Class 1 has
a higher threshold, reflecting that the more tolerant class captures
more available slots in FCFS.}{fig:balk-deep-threshold-gap}

%% ============================================================
\section{Regression Screen}
\label{sec:regression}
%% ============================================================

The regression screen tests whether the main visual findings hold when all
assumptions vary simultaneously. For each of 240 randomized FCFS parameter
settings, the simulation is run with two seeds and the outputs averaged.
Twelve features are constructed from the sampled inputs: total arrival rate,
Class~1 arrival share, average levels, and Class~1 minus Class~2 gaps for
balking step, balking threshold, no-show step, no-show threshold, and
cancellation probability. Four targets are modeled separately.

All variables are standardized before fitting ordinary least squares.
Significance is assessed using HC3 robust standard errors; only coefficients
with robust p-values below 0.05 are reported. Full-sample $R^2$ ranges from
0.639 (class access advantage) to 0.799 (overall served rate), and
out-of-sample $R^2$ from 0.545 (utilization) to 0.719 (overall served rate),
indicating that the linear screen explains meaningful variation in the
simulation design.

These are simulation associations within the randomized design, not
causal estimates from real appointment data.

\begin{table}[H]
\centering
\small
\renewcommand{\arraystretch}{1.08}
\begin{tabular}{p{0.15\textwidth}p{0.26\textwidth}rp{0.20\textwidth}r}
\toprule
Target metric & Significant feature & Std.\ coef.\ & 95\% HC3 CI & HC3 $p$ \\
\midrule
Utilization & Average no-show threshold & $+0.512$ & [0.433, 0.591] & $<0.001$ \\
Utilization & Average no-show step & $-0.423$ & [$-0.501$, $-0.346$] & $<0.001$ \\
Utilization & Average cancellation probability & $+0.320$ & [0.247, 0.393] & $<0.001$ \\
Utilization & Total arrival rate & $-0.109$ & [$-0.188$, $-0.031$] & 0.006 \\
Utilization & Average balking threshold & $-0.086$ & [$-0.162$, $-0.009$] & 0.028 \\
\midrule
Overall served rate & Total arrival rate & $-0.774$ & [$-0.855$, $-0.694$] & $<0.001$ \\
Overall served rate & Average no-show threshold & $+0.251$ & [0.193, 0.308] & $<0.001$ \\
Overall served rate & Average no-show step & $-0.175$ & [$-0.232$, $-0.119$] & $<0.001$ \\
Overall served rate & Average cancellation probability & $+0.117$ & [0.058, 0.176] & $<0.001$ \\
\midrule
Offered wait & Total arrival rate & $+0.576$ & [0.499, 0.653] & $<0.001$ \\
Offered wait & Average cancellation probability & $-0.441$ & [$-0.506$, $-0.376$] & $<0.001$ \\
Offered wait & Average balking threshold & $+0.278$ & [0.202, 0.353] & $<0.001$ \\
Offered wait & Average balking step & $-0.206$ & [$-0.291$, $-0.122$] & $<0.001$ \\
\midrule
Class gap & Cancellation probability gap & $-0.452$ & [$-0.558$, $-0.347$] & $<0.001$ \\
Class gap & Balking threshold gap & $+0.429$ & [0.331, 0.528] & $<0.001$ \\
Class gap & Balking step gap & $-0.349$ & [$-0.431$, $-0.266$] & $<0.001$ \\
Class gap & No-show threshold gap & $+0.268$ & [0.179, 0.357] & $<0.001$ \\
Class gap & No-show step gap & $-0.211$ & [$-0.286$, $-0.136$] & $<0.001$ \\
Class gap & Class~1 arrival share & $-0.120$ & [$-0.209$, $-0.030$] & 0.009 \\
\bottomrule
\end{tabular}
\caption{Significant regression drivers ($p < 0.05$, HC3 robust). Gap
features are Class~1 value minus Class~2 value. Coefficients are in
standard deviation units.}
\label{tab:regression}
\end{table}

\mainfigure{regression_significant_standardized_coefficients.png}{%
Regression coefficients confirm the metric-by-metric reading: no-show
behavior dominates utilization, demand pressure dominates served rate
and offered wait, and class behavioral gaps dominate the class
served-rate gap. Error bars are 95\% HC3 robust confidence
intervals.}{fig:metric-regression-full}

%% ============================================================
\section{Key Findings}
%% ============================================================

\begin{enumerate}

\item \textbf{High utilization can coexist with poor access.}
At baseline, utilization is 0.839 but overall served rate is 0.269.
The system fills available slots while failing to serve the majority of
arriving patients. Utilization should not be used as a standalone
performance metric in this scheduling context.

\item \textbf{Overall served rate is the cleanest access metric.}
It uses all arrivals as the denominator and counts only completed visits.
It is not affected by balking selection or the distinction between offered
and accepted delays.

\item \textbf{Mean offered wait is preferred over mean accepted wait
for patient-facing delay.}
Accepted wait is conditional on patients accepting the offered appointment.
It falls mechanically when patients with long offered delays balk, giving
the misleading impression that congestion has improved. Offered wait includes
those patients and better reflects the delay that the system actually imposes.

\item \textbf{Balking can mask congestion.}
The deep-dive analysis shows that increasing Class~1 balking probability
reduces Class~1 accepted wait while simultaneously reducing Class~1 access
and potentially reducing overall offered wait. Interpreting lower accepted
wait as a service improvement in this setting is incorrect.

\item \textbf{Demand pressure is the dominant driver of access; behavioral
parameters dominate class equity.}
The regression screen assigns the largest absolute standardized coefficient
for overall served rate to total arrival rate ($|\beta| = 0.774$). For
class served-rate gap, the largest coefficient belongs to the cancellation
probability gap ($|\beta| = 0.452$), followed by balking threshold gap and
balking step gap.

\item \textbf{Class disparities arise from behavioral asymmetries, not
from class labels.}
In the symmetric baseline, the class gap is near zero (0.001). Gap magnitude
and direction track differences in cancellation probability, balking
threshold, balking step, and no-show parameters. Equity-relevant
interventions should therefore target behavioral parameter differences
rather than class membership per se.

\item \textbf{Regression associations are simulation-internal and should
not be interpreted as real-world causal effects.}
The randomized screen shows stable linear associations within the
simulation design. Real-world behavioral parameters are not observed;
they are assumed. Translating these findings to clinic policy requires
empirical calibration.

\end{enumerate}

%% ============================================================
\section{Limitations and Next Steps}
%% ============================================================

\subsection{Limitations}

\begin{enumerate}

\item \textbf{Stylized simulation design.}
The model uses threshold-type behavioral rules with sharp transition
probabilities. Real patients likely have smoother, heterogeneous delay
tolerances. The sharp thresholds amplify sensitivity around the threshold
values in ways that may not appear in observed data.

\item \textbf{Assumed behavioral parameters.}
Balking probabilities, no-show probabilities, and cancellation rates are
set by assumption, not estimated from patient data. The sensitivity analysis
covers a wide parameter range precisely because the true values are unknown.
Any policy inference requires calibration against clinic records.

\item \textbf{Limited replication.}
Most sensitivity grid points use one to three random seeds. Metrics with
high seed-to-seed variance (particularly the class gap at small
parameter differences) may not have converged. The regression screen averages
only two seeds per setting.

\item \textbf{Single scheduling discipline.}
The entire analysis uses FCFS. FCFS is a reasonable baseline but is not
the only feasible policy. Priority-based, panel-based, or advanced-access
scheduling may have qualitatively different utilization--access tradeoffs.

\item \textbf{No provider-side dynamics.}
The model holds slots per day fixed. Real clinics adjust templates, add
sessions, or use overbooking. These supply-side responses are not modeled.

\end{enumerate}

\subsection{Next Steps}

\begin{enumerate}

\item \textbf{Calibrate behavioral parameters from clinic data.}
Estimate patient-specific balking and no-show rates using appointment
history and no-show records. Use these estimates to anchor the
sensitivity grid around empirically plausible ranges.

\item \textbf{Evaluate alternative scheduling disciplines.}
Implement and compare at least one priority-based policy (e.g., by
clinical urgency or social risk) to assess whether FCFS produces
systematically worse access for higher-need classes.

\item \textbf{Increase simulation replications.}
Run at least 10--20 seeds per grid point for the primary metrics and
at least 30--50 for the class gap, which is more variable. Report
confidence intervals on metric estimates.

\item \textbf{Model supply-side feedback.}
Introduce session-level cancellation and overbooking rules to test
how the clinic-side response to no-shows affects the
utilization--access gap.

\item \textbf{Extend to provider panel and network effects.}
If patient panels are partitioned across providers, FCFS within a
panel may produce different access outcomes than pooled FCFS. Modeling
panel structure would allow equity analysis at the provider level.

\end{enumerate}

%% ============================================================
\appendix
\section{Appendix: Supplementary Sensitivity Plots}
%% ============================================================

These figures are retained for traceability. They support the main text
but are not required for the primary findings.

\subsection{No-Show Sensitivity}

\mainfigure{no_show_step_interaction_heatmap.png}{%
No-show step heatmap. Left: utilization falls when both classes have
high no-show risk. Right: class gap moves against the class with
higher no-show step.}{fig:app-noshow-step}

\mainfigure{no_show_threshold_interaction_heatmap.png}{%
No-show threshold heatmap. A higher threshold delays the onset of
high no-show risk; utilization and access improve for the affected
class.}{fig:app-noshow-threshold}

\subsection{Balking Sensitivity}

\mainfigure{balking_step_interaction_heatmap.png}{%
Balking step heatmap. Left: overall served rate falls as either class
increases its high-balking probability. Right: class gap moves
against the class with higher balking step.}{fig:app-balk-step}

\mainfigure{balking_threshold_interaction_heatmap.png}{%
Balking threshold heatmap. Higher threshold allows longer waits
before high-balking applies; the more tolerant class captures more
served slots under FCFS.}{fig:app-balk-threshold}

\mainfigure{balking_step_balking_rate_heatmap.png}{%
Balking rate under balking-step changes. The right panel shows the
Class~1 minus Class~2 balking-rate gap.}{fig:app-balk-rate-step}

\mainfigure{balking_threshold_balking_rate_heatmap.png}{%
Balking rate under balking-threshold changes. The more patient class
has the lower balking rate and the higher served rate.
}{fig:app-balk-rate-threshold}

\subsection{Demand and Cancellation}

\mainfigure{fcfs_arrival_outcome_decomposition.png}{%
Outcome decomposition as total demand rises. At high demand, many
arrivals receive no offer because the FCFS horizon fills; at moderate
demand, balking and after-booking losses dominate.
}{fig:app-outcome-decomp}

\mainfigure{arrival_mix_interaction_heatmap.png}{%
Arrival rate and class-mix heatmap. Left: overall served rate.
Right: mean offered wait. Higher total demand worsens both.
}{fig:app-arrival-mix}

\mainfigure{arrival_rate_slice_wait.png}{%
Mean offered wait as total arrivals change at a 50/50 class mix.
Offered wait rises monotonically with demand.
}{fig:app-arrival-wait}

\mainfigure{cancellation_probability_interaction_heatmap.png}{%
Cancellation heatmap. Left: overall served rate. Right: Class~1
served-rate gap. Higher cancellation for one class reduces that
class's access.}{fig:app-cancel}

\subsection{Common Threshold Grids}

\mainfigure{balking_threshold_jump_heatmaps.png}{%
Shared balking threshold and jump level for both classes. These are
symmetric aggregate tests, not class advantage tests.
}{fig:app-balk-common}

\mainfigure{no_show_threshold_jump_heatmaps.png}{%
Shared no-show threshold and jump level for both classes.
}{fig:app-noshow-common}

\subsection{Full Diagnostic Plots}

\mainfigure{fcfs_capacity_stress_curves.png}{%
Full FCFS stress diagnostic plot across the wider arrival range.
}{fig:app-fcfs-full}

\subsection{Balking Deep Dive: Additional Plots}

\mainfigure{../../../../../outputs/class1_balking/figures/average_utilization_aggregate.png}{%
Class~1 high-balking-probability sweep: aggregate utilization remains
nearly flat, indicating that Class~2 demand fills the slots vacated by
balking Class~1 patients.}{fig:app-balk-deep-util}

\mainfigure{../../../../../outputs/class1_balking/figures/balking_rate_by_class.png}{%
Class~1 high-balking-probability sweep: balking-rate confirmation that
the mechanism is behavioral rejection, not capacity change.
}{fig:app-balk-deep-rate}

\mainfigure{../../../../../outputs/class1_balking_threshold/figures/average_utilization_aggregate.png}{%
Class~1 balking-threshold sweep: aggregate utilization is largely
flat, confirming that threshold changes redistribute service without
large aggregate effects.}{fig:app-balk-thresh-util}

\mainfigure{../../../../../outputs/class1_balking_threshold/figures/mean_accepted_delay_by_class.png}{%
Class~1 balking-threshold sweep: Class~1 accepted delay rises with
the threshold because patients tolerate longer waits before the high-
balking probability applies.}{fig:app-balk-thresh-accepted}

\mainfigure{../../../../../outputs/class1_balking_threshold/figures/mean_offered_delay_by_class.png}{%
Class~1 balking-threshold sweep: mean offered delay by class.
}{fig:app-balk-thresh-offered}

%% ============================================================
\section{Appendix: Exact Metric Definitions}
%% ============================================================

Let $i$ be patient class, $D$ be measured service days, $S$ be slots per
day, $A_i$ be arrivals, $B_i$ be accepted bookings, $K_i$ be balked patients,
$C_i$ be cancellations, $H_i$ be no-shows, $V_i$ be completed visits, and
$O_i = B_i + K_i$ be offered patients. Let $\bar{\tau}^{\text{offered}}_i$
be class $i$ mean offered booking delay.

\begin{align*}
\texttt{percent\_serviced}_i &= \frac{V_i}{A_i}, \\
\texttt{overall\_percent\_serviced} &= \frac{\sum_i V_i}{\sum_i A_i}, \\
\texttt{mean\_offered\_booking\_delay} &= \frac{\sum_i \text{total offered delay}_i}{\sum_i O_i}, \\
\texttt{mean\_accepted\_booking\_delay} &= \frac{\sum_i \text{total accepted delay}_i}{\sum_i B_i}, \\
\texttt{average\_utilization} &= \frac{1}{D}\sum_d \frac{\text{served slots}_d}{S}, \\
\texttt{balking\_rate} &= \frac{\sum_i K_i}{\sum_i O_i}, \\
\Delta_{\text{access}} &= \texttt{percent\_serviced}_1 - \texttt{percent\_serviced}_2, \\
\Delta_{\text{delay}} &= \bar{\tau}^{\text{offered}}_2 - \bar{\tau}^{\text{offered}}_1, \\
\Delta_{\text{balking}} &= \frac{K_1}{O_1} - \frac{K_2}{O_2}.
\end{align*}

\noindent Positive $\Delta_{\text{access}}$ means Class~1 has the higher
served rate. Positive $\Delta_{\text{delay}}$ means Class~1 has the lower
mean offered wait (Class~2 waits longer). Positive $\Delta_{\text{balking}}$
means Class~1 rejects offers more often.

\noindent\textbf{Measurement semantics note.} Class-level patient metrics
are tracked for patients who arrive during the measurement window. If
\texttt{cooldown\_days} $<$ \texttt{horizon\_days} $- 1$, some late
measurement-window bookings may remain unresolved when the simulation ends;
these patients are currently not counted as served, canceled, or no-show.
\texttt{total\_value} is service-day based and is not cohort-aligned with
the tracked class metrics.

\end{document}
